\section{Open source}

Open source označuje software, jehož zdrojový kód je dostupný a jehož distribuce je upravena licencí splňující kritéria Open Source Definition (OSD). Důležité je, že open source neznamená pouze „přístup ke kódu“, ale zejména vymezení práv a povinností při používání, úpravách a šíření softwaru. \parencite{OpenSourceDefinition,OpenSource2024}

\subsection{Open source licence}
Podmínky použití a šíření open source softwaru jsou dány licencí. Open Source Initiative (OSI) udržuje definici OSD a zároveň zveřejňuje seznam licencí, které jsou s OSD v souladu (tzv.~OSI Approved). \parencite{OpenSourceDefinition,OSILicenses}

\subsection{Copyleft a permisivní licence}
Významnou skupinou open-source licencí jsou copyleft licence, které při šíření odvozených děl typicky vyžadují zachování stejných (či kompatibilních) licenčních podmínek. Oproti tomu permisivní licence ponechávají příjemci větší volnost v dalším šíření a v licencování odvozených děl. \parencite{Copyleft2026}

\subsection{Relevance k této práci}
V kontextu této práce je open source relevantní zejména z pohledu transparentnosti a znovupoužitelnosti implementovaného řešení. Zvolená licence ovlivňuje, za jakých podmínek může být aplikace dále používána, upravována a případně rozvíjena dalšími uživateli.


